\section{Verificación por Simulación SPICE y Parámetros de Calidad}
\label{sec:simulaciones}

\subsection{Verificación de la Regulación de Línea para $V_{in} = \SI{25}{\volt}$ y $V_{in} = \SI{10}{\volt}$}
Con el objeto de verificar experimentalmente en entorno de simulación el desempeño del diseño, se modeló el sistema completo en \texttt{ngspice} (\texttt{sim/fuente\_7805.cir}) y \texttt{LTspice} (\texttt{sim/fuente\_7805.asc}). El macromodelo del regulador integrado \texttt{LM7805C} reproduce fielmente las características eléctricas de la hoja de datos de la Figura 6-59 de \cite{gonzalez_p3} (tensión nominal \SI{5,000}{\volt} a $V_I = \SI{10}{\volt}$ e $I_O = \SI{500}{\milli\ampere}$, regulación de línea de \SI{1,0}{\milli\volt} entre \SI{8}{\volt} y \SI{12}{\volt} y de \SI{3,2}{\milli\volt} entre \SI{7}{\volt} y \SI{25}{\volt}, tensión de \emph{dropout} de \SI{2,0}{\volt}, impedancia de salida $r_o = \SI{0,017}{\ohm}$, corriente quiescente $I_d = \SI{4,2}{\milli\ampere}$ y limitación de cortocircuito en \SI{1,5}{\ampere}).

Se ejecutó un barrido continuo de la tensión de entrada primaria $V_{in}$ desde \SI{0}{\volt} hasta \SI{30}{\volt} con carga nominal $R_L = \SI{10,0}{\ohm}$ ($I_L = \SI{500}{\milli\ampere}$) comparando tres configuraciones:
\begin{itemize}
    \item \textbf{\texttt{Circuito A} (Preregulador de rango extendido $\SI{10}{\volt}\text{--}\SI{25}{\volt}$):} $Q_1 = \texttt{TIP31C}$, $D_Z = \texttt{1N4740A}$ (\SI{10}{\volt}, \SI{1}{\watt}), $R_Z = \SI{180}{\ohm}$.
    \item \textbf{\texttt{Circuito B} (Preregulador dimensionado en \SI{25}{\volt} según Secc. 6-9-5h):} $Q_1 = \texttt{TIP31C}$, $D_Z = \texttt{1N5240B}$ (\SI{10}{\volt}, \SI{500}{\milli\watt}), $R_Z = \SI{430}{\ohm}$.
    \item \textbf{\texttt{Circuito C} (Referencia sin etapa prereguladora):} Regulador \texttt{LM7805C} conectado de forma directa a la fuente primaria $V_{in}$.
\end{itemize}

\begin{figure*}[t]
\centering
\begin{tikzpicture}
\begin{axis}[
    name=plot1,
    width=0.46\textwidth,
    height=6.2cm,
    xlabel={Tensión de Entrada Primaria $V_{in}$ [\si{\volt}]},
    ylabel={Tensiones de Nodo [\si{\volt}]},
    xmin=0, xmax=30,
    ymin=0, ymax=14,
    grid=both,
    grid style={line width=.1pt, draw=gray!25},
    major grid style={line width=.2pt, draw=gray!50},
    legend pos=south east,
    legend style={font=\scriptsize, fill=white, fill opacity=0.92, draw=gray!60},
    tick label style={font=\footnotesize},
    label style={font=\footnotesize\bfseries},
    title style={font=\footnotesize\bfseries, align=center},
    title={(a) Acción de la Etapa Prereguladora sobre $V_{i,\text{reg}}$ y $V_O$}
]
    % Banda sombreada 8 V - 12 V
    \addplot[name path=top12, draw=none, forget plot] coordinates {(0,12) (30,12)};
    \addplot[name path=bot8, draw=none, forget plot] coordinates {(0,8) (30,8)};
    \addplot[green!25, opacity=0.45] fill between[of=top12 and bot8];
    \addlegendentry{Rango objetivo $[\SI{8}{\volt}, \SI{12}{\volt}]$}

    % Curvas de tensión
    \addplot[color=orange!85!black, thick, dashed] table[x=Vin, y=VzA] {sim/regulacion_linea.txt};
    \addlegendentry{$V_Z$ Base $Q_1$ (\texttt{Circ. A})}

    \addplot[color=blue!80!black, ultra thick] table[x=Vin, y=ViregA] {sim/regulacion_linea.txt};
    \addlegendentry{$V_{i,\text{reg}}$ Entrada \texttt{7805} (\texttt{Circ. A})}

    \addplot[color=teal!80!black, thick, dashdotted] table[x=Vin, y=ViregB] {sim/regulacion_linea.txt};
    \addlegendentry{$V_{i,\text{reg}}$ Entrada \texttt{7805} (\texttt{Circ. B})}

    \addplot[color=red!85!black, ultra thick] table[x=Vin, y=VoA] {sim/regulacion_linea.txt};
    \addlegendentry{Salida regulada $V_O$ (\SI{5,00}{\volt})}

    % Puntos destacados en Vin = 10 V y Vin = 25 V
    \addplot[only marks, mark=*, mark size=2.4pt, color=blue!85!black, forget plot] coordinates {
        (10.0, 8.056)
        (25.0, 9.543)
    };
    \addplot[only marks, mark=*, mark size=2.4pt, color=red!85!black, forget plot] coordinates {
        (10.0, 4.997)
        (25.0, 4.999)
    };
    \node[anchor=south east, font=\scriptsize\bfseries, text=blue!80!black] at (axis cs:10.5, 8.25) {$(\SI{10}{\volt}; \SI{8,06}{\volt})$};
    \node[anchor=north, font=\scriptsize\bfseries, text=blue!80!black] at (axis cs:24.2, 9.15) {$(\SI{25}{\volt}; \SI{9,54}{\volt})$};
\end{axis}

\begin{axis}[
    name=plot2,
    at={(plot1.outer north east)},
    anchor=outer north west,
    xshift=0.75cm,
    width=0.46\textwidth,
    height=6.2cm,
    xlabel={Tensión de Entrada Primaria $V_{in}$ [\si{\volt}]},
    ylabel={Variación $\Delta V_O = V_O - V_{O(\SI{10}{\volt})}$ [\si{\milli\volt}]},
    xmin=9.5, xmax=30,
    ymin=-1.5, ymax=10.5,
    restrict x to domain=9.5:30,
    restrict y to domain=-10:20,
    grid=both,
    grid style={line width=.1pt, draw=gray!25},
    major grid style={line width=.2pt, draw=gray!50},
    legend pos=north west,
    legend style={font=\scriptsize, fill=white, fill opacity=0.92, draw=gray!60},
    tick label style={font=\footnotesize},
    label style={font=\footnotesize\bfseries},
    title style={font=\footnotesize\bfseries, align=center},
    title={(b) Comparación de Regulación de Línea (Con vs. Sin Preregulador)}
]
    % Sin preregulador
    \addplot[color=purple!80!black, ultra thick, dashed] table[x=Vin, y=dVoDirecto_mV] {sim/regulacion_linea.txt};
    \addlegendentry{Sin Preregulador (\texttt{LM7805C} directo)}

    % Con preregulador
    \addplot[color=blue!85!black, ultra thick] table[x=Vin, y=dVoA_mV] {sim/regulacion_linea.txt};
    \addlegendentry{Con Preregulador (\texttt{Circuito A})}

    % Puntos en Vin = 10 V y Vin = 25 V
    \addplot[only marks, mark=*, mark size=2.4pt, color=purple!80!black, forget plot] coordinates {
        (10.0, 0.0)
        (25.0, 6.79)
    };
    \addplot[only marks, mark=*, mark size=2.4pt, color=blue!85!black, forget plot] coordinates {
        (10.0, 0.0)
        (25.0, 1.70)
    };

    \node[anchor=north east, font=\scriptsize\bfseries, text=purple!80!black] at (axis cs:29.5, 6.45) {$\Delta V_{O(\text{sin})} = \SI{6,79}{\milli\volt}$ (a \SI{25}{\volt})};
    \node[anchor=north east, font=\scriptsize\bfseries, text=blue!85!black] at (axis cs:29.5, 1.45) {$\Delta V_{O(\text{con})} = \SI{1,70}{\milli\volt}$ (\SI{0,28}{\milli\volt} de \SI{12}{\volt} a \SI{25}{\volt})};
\end{axis}
\end{tikzpicture}
\caption{Resultados de simulación SPICE del barrido de línea: (a) estabilización de la tensión intermedia $V_{i,\text{reg}}$ dentro de la banda $[\SI{8}{\volt}, \SI{12}{\volt}]$ tanto en $V_{in} = \SI{10}{\volt}$ como en $V_{in} = \SI{25}{\volt}$; (b) mejora de la regulación de línea respecto del regulador conectado directamente a la fuente primaria.}
\label{fig:simulacion_linea}
\end{figure*}

En la Figura~\ref{fig:simulacion_linea} y en la Tabla~\ref{tab:resultados_linea} se detallan los resultados obtenidos de la simulación:
\begin{enumerate}
    \item \textbf{Cumplimiento estricto del rango $V_{i,\text{reg}} \in [\SI{8}{\volt}, \SI{12}{\volt}]$:} Como se aprecia en la Figura~\ref{fig:simulacion_linea}(a), cuando la fuente primaria opera en su valor nominal de \textbf{$V_{in} = \SI{25,0}{\volt}$}, la etapa prereguladora fija la tensión de entrada al \texttt{LM7805C} en \textbf{$V_{i,\text{reg}} = \SI{9,543}{\volt}$} (\texttt{Circuito A}) y \textbf{$V_{i,\text{reg}} = \SI{9,330}{\volt}$} (\texttt{Circuito B}, coincidiendo con los \SI{9,388}{\volt} del ejemplo de cátedra en la pág. 17 de \cite{gonzalez_p3}). Al reducir la fuente primaria hasta \textbf{$V_{in} = \SI{10,0}{\volt}$}, el \texttt{Circuito A} mantiene \textbf{$V_{i,\text{reg}} = \SI{8,056}{\volt}$}, cumpliendo que durante toda la excursión de \SI{15}{\volt} en la línea primaria ($\SI{10}{\volt} \le V_{in} \le \SI{25}{\volt}$) el integrado trabaja ininterrumpidamente dentro de la ventana óptima $[\SI{8,0}{\volt}, \SI{12,0}{\volt}]$.
    \item \textbf{Comparación de regulación de línea en $V_{in} = \SI{25}{\volt}$ vs. $V_{in} = \SI{10}{\volt}$:}
    \begin{itemize}
        \item \emph{Excursión completa de $V_{in} = \SI{10,0}{\volt}$ a $V_{in} = \SI{25,0}{\volt}$ ($\Delta V_{in} = \SI{15,0}{\volt}$):} Sin etapa prereguladora (\texttt{Circuito C}), el terminal \texttt{IN} del \texttt{LM7805C} sufre los \SI{15,0}{\volt} de variación, produciendo un cambio en la salida de $\Delta V_{O(\text{sin})} = \SI{6,790}{\milli\volt}$ ($R_{\text{red}} = \SI{0,1358}{\percent}$). Con la etapa prereguladora (\texttt{Circuito A}), la variación en $V_{i,\text{reg}}$ se acota a sólo $\Delta V_{i,\text{reg}} = \SI{9,543}{\volt} - \SI{8,056}{\volt} = \SI{1,487}{\volt}$, reduciendo la variación de salida a \textbf{$\Delta V_{O(\text{con})} = \SI{1,699}{\milli\volt}$} ($R_{\text{red}} = \SI{0,0340}{\percent}$, mejora de $4\times$). Más aún, dentro de la región de ruptura plena del Zener ($\SI{12,0}{\volt} \le V_{in} \le \SI{25,0}{\volt}$, $\Delta V_{in} = \SI{13,0}{\volt}$), la variación de salida cae a apenas \textbf{$\Delta V_O = \SI{0,277}{\milli\volt}$} (\texttt{Circuito A}) y \textbf{$\SI{0,938}{\milli\volt}$} (\texttt{Circuito B}).
        \item \emph{Fluctuaciones normalizadas de red del $\pm \SI{10}{\percent}$ en torno a $V_{in} = \SI{25,0}{\volt}$ ($\SI{22,5}{\volt}$ a $\SI{27,5}{\volt}$) y en torno a $V_{in} = \SI{10,0}{\volt}$ ($\SI{9,0}{\volt}$ a $\SI{11,0}{\volt}$):} Cuando el \texttt{LM7805C} trabaja sin preregulador en $V_{in} = \SI{25,0}{\volt} \pm \SI{10}{\percent}$, presenta $\Delta V_O = \SI{2,792}{\milli\volt}$; cuando trabaja directamente en $V_{in} = \SI{10,0}{\volt} \pm \SI{10}{\percent}$ (dentro de la zona óptima de \SI{8}{\volt} a \SI{12}{\volt}), su variación baja a $\Delta V_O = \SI{1,228}{\milli\volt}$. Al intercalar la etapa prereguladora diseñada operando desde la línea de $\SI{25,0}{\volt} \pm \SI{10}{\percent}$, la variación a la salida se reduce a tan sólo \textbf{$\Delta V_O = \SI{0,084}{\milli\volt}$ (\SI{84,1}{\micro\volt})} en el \texttt{Circuito A} y \textbf{$\Delta V_O = \SI{0,047}{\milli\volt}$ (\SI{46,5}{\micro\volt})} en el \texttt{Circuito B}, lo que representa una \textbf{mejora de $33\times$ a $60\times$} en la regulación de línea.
    \end{itemize}
\end{enumerate}

\begin{table}[htbp]
\centering
\caption{Resultados de simulación SPICE para la regulación de línea ($I_L = \SI{500}{\milli\ampere}$, $R_L = \SI{10,0}{\ohm}$).}
\label{tab:resultados_linea}
\resizebox{\columnwidth}{!}{%
\begin{tabular}{lccccc}
\toprule
\textbf{Condición de Entrada $V_{in}$} & \textbf{$V_Z$ [\si{\volt}]} & \textbf{$V_{i,\text{reg}}$ [\si{\volt}]} & \textbf{$V_{O(\text{con})}$ [\si{\volt}]} & \textbf{$V_{O(\text{sin})}$ [\si{\volt}]} & \textbf{$P_{D(\text{7805})}$ Con / Sin [\si{\watt}]} \\
\midrule
$V_{in} = \SI{10,0}{\volt}$ (\texttt{Circ. A}) & \SI{8,771}{} & \textbf{\SI{8,056}{}} & \SI{4,99704}{} & \SI{4,99903}{} & \SI{1,56}{} / \SI{2,54}{} \\
$V_{in} = \SI{12,0}{\volt}$ (\texttt{Circ. A / B}) & \SI{9,898}{} & \textbf{\SI{9,184}{}} & \SI{4,99846}{} & \SI{4,99995}{} & \SI{2,13}{} / \SI{3,55}{} \\
$V_{in} = \SI{22,5}{\volt}$ ($\SI{25}{\volt} - \SI{10}{\percent}$) & \SI{10,193}{} & \textbf{\SI{9,483}{}} & \SI{4,99870}{} & \SI{5,00449}{} & \SI{2,28}{} / \SI{8,85}{} \\
$V_{in} = \SI{25,0}{\volt}$ (Nominal \texttt{Circ. A}) & \SI{10,253}{} & \textbf{\SI{9,543}{}} & \SI{4,99874}{} & \SI{5,00582}{} & \textbf{\SI{2,31}{} / \SI{10,11}{}} \\
$V_{in} = \SI{25,0}{\volt}$ (Nominal \texttt{Circ. B}) & \SI{10,040}{} & \textbf{\SI{9,330}{}} & \SI{4,99858}{} & \SI{5,00582}{} & \textbf{\SI{2,20}{} / \SI{10,11}{}} \\
$V_{in} = \SI{27,5}{\volt}$ ($\SI{25}{\volt} + \SI{10}{\percent}$) & \SI{10,312}{} & \textbf{\SI{9,603}{}} & \SI{4,99878}{} & \SI{5,00728}{} & \SI{2,34}{} / \SI{11,37}{} \\
\midrule
\textbf{Variación $\SI{10}{\volt} \to \SI{25}{\volt}$} & $\Delta V_Z = \SI{1,48}{\volt}$ & $\Delta V_{i,\text{reg}} = \SI{1,49}{\volt}$ & $\mathbf{\Delta V_O = \SI{1,70}{\milli\volt}}$ & $\Delta V_O = \SI{6,79}{\milli\volt}$ & Reducción del \SI{77,2}{\percent} \\
\textbf{Variación $\SI{12}{\volt} \to \SI{25}{\volt}$} & $\Delta V_Z = \SI{0,35}{\volt}$ & $\Delta V_{i,\text{reg}} = \SI{0,36}{\volt}$ & $\mathbf{\Delta V_O = \SI{0,28}{\milli\volt}}$ & $\Delta V_O = \SI{5,87}{\milli\volt}$ & Mejora de $21\times$ \\
\textbf{Variación $\SI{25}{\volt} \pm \SI{10}{\percent}$} & $\Delta V_Z = \SI{0,12}{\volt}$ & $\Delta V_{i,\text{reg}} = \SI{0,12}{\volt}$ & $\mathbf{\Delta V_O = \SI{0,084}{\milli\volt}}$ & $\Delta V_O = \SI{2,79}{\milli\volt}$ & Mejora de $33\times$ (\texttt{B}: $60\times$) \\
\bottomrule
\end{tabular}%
}
\end{table}

\subsection{Cálculo y Simulación de los Parámetros de Calidad de la Fuente}
Siguiendo estrictamente las definiciones de la Sección 6-6-3 del material de cátedra \cite{gonzalez_p2}, se determinan a continuación todos los factores de mérito de la fuente diseñada.

\subsubsection{Regulación de Red ($R_{\text{red}}$)}
Definida por la Ecuación 6-41 de \cite{gonzalez_p2} como la variación porcentual de la tensión de salida sobre la carga frente a variaciones entre los límites mínimo y máximo de la tensión de entrada, manteniendo constantes la corriente de carga ($I_L = \SI{500}{\milli\ampere}$) y la temperatura ($T_J = \SI{25}{\celsius}$):
\begin{equation}
    R_{\text{red}} = \frac{V_{L(\text{red máx})} - V_{L(\text{red mín})}}{V_{L(\text{red mín})}} \times 100 \quad [\%]
    \label{eq:rred_def}
\end{equation}
\begin{itemize}
    \item \textbf{Ante excursión extrema de línea de $V_{in} = \SI{10,0}{\volt}$ a $V_{in} = \SI{25,0}{\volt}$ ($\Delta V_O = \SI{1,70}{\milli\volt}$):}
    \begin{equation}
        R_{\text{red}(\SI{10}{\volt}\text{--}\SI{25}{\volt})} = \frac{\SI{4,99874}{\volt} - \SI{4,99704}{\volt}}{\SI{4,99704}{\volt}} \times 100 = \mathbf{\SI{0,0340}{\percent}}
    \end{equation}
    \item \textbf{Ante variación de línea en régimen Zener activo de $\SI{12,0}{\volt}$ a $\SI{25,0}{\volt}$ ($\Delta V_O = \SI{0,28}{\milli\volt}$):}
    \begin{equation}
        R_{\text{red}(\SI{12}{\volt}\text{--}\SI{25}{\volt})} = \frac{\SI{4,99874}{\volt} - \SI{4,99846}{\volt}}{\SI{4,99846}{\volt}} \times 100 = \mathbf{\SI{0,0055}{\percent}}
    \end{equation}
    \item \textbf{Ante fluctuación normalizada de red de $\pm \SI{10}{\percent}$ en torno a \SI{25,0}{\volt} ($\SI{22,5}{\volt}$ a $\SI{27,5}{\volt}$, $\Delta V_O = \SI{84,1}{\micro\volt}$ en \texttt{Circ. A} y \SI{46,5}{\micro\volt} / \SI{0,00093}{\percent} en \texttt{Circ. B}):}
    \begin{equation}
        R_{\text{red}(\pm 10\%)} = \frac{\SI{4,99878}{\volt} - \SI{4,99870}{\volt}}{\SI{4,99870}{\volt}} \times 100 = \mathbf{\SI{0,00168}{\percent}}
    \end{equation}
\end{itemize}

\subsubsection{Regulación de Carga ($R_{\text{car}}$) y Resistencia de Salida ($R_{\text{sal}}$)}
La regulación de carga (Ec. 6-42 de \cite{gonzalez_p2}) cuantifica la caída porcentual de la tensión de salida cuando la corriente demandada varía desde carga mínima ($I_{L,\min}$) hasta carga máxima ($I_{L,\max}$), a tensión de entrada constante ($V_{in} = \SI{25,0}{\volt}$):
\begin{equation}
    R_{\text{car}} = \frac{V_{L(\text{car mín})} - V_{L(\text{car máx})}}{V_{L(\text{car máx})}} \times 100 \quad [\%]
    \label{eq:rcar_def}
\end{equation}
A partir del barrido de carga simulado en \texttt{ngspice} (\texttt{sim/regulacion\_carga.txt}, graficado en la Figura~\ref{fig:simulacion_carga_ripple}(a)):
\begin{itemize}
    \item Para el rango de $I_{L,\min} = \SI{5}{\milli\ampere}$ ($V_L = \SI{5,00725}{\volt}$) hasta la corriente nominal $I_{L,\text{nom}} = \SI{500}{\milli\ampere}$ ($V_L = \SI{4,99874}{\volt}$, $\Delta V_O = \SI{8,51}{\milli\volt}$):
    \begin{equation}
        R_{\text{car}(\SI{0,5}{\ampere})} = \frac{\SI{5,00725}{\volt} - \SI{4,99874}{\volt}}{\SI{4,99874}{\volt}} \times 100 = \mathbf{\SI{0,170}{\percent}}
    \end{equation}
    \item Para el rango completo hasta la corriente máxima continua $I_{L,\max} = \SI{1,00}{\ampere}$ ($V_L = \SI{4,99010}{\volt}$, $\Delta V_O = \SI{17,15}{\milli\volt}$, hoja de datos: \SI{15}{\milli\volt} típ.):
    \begin{equation}
        R_{\text{car}(\SI{1,0}{\ampere})} = \frac{\SI{5,00725}{\volt} - \SI{4,99010}{\volt}}{\SI{4,99010}{\volt}} \times 100 = \mathbf{\SI{0,344}{\percent}}
    \end{equation}
\end{itemize}
La \textbf{resistencia estática y dinámica de salida} ($R_{\text{sal}}$, Ec. 6-43 de \cite{gonzalez_p2}) resulta:
\begin{equation}
    R_{\text{sal}} = \frac{V_{L(\text{car mín})} - V_{L(\text{car máx})}}{I_{L(\text{car máx})} - I_{L(\text{car mín})}} = \frac{\SI{17,15}{\milli\volt}}{\SI{0,995}{\ampere}} = \mathbf{\SI{0,0172}{\ohm}} = \SI{17,2}{\milli\ohm}
    \label{eq:rsal_calc}
\end{equation}
valor idéntico a la impedancia dinámica de hoja de datos $r_o = \SI{0,017}{\ohm}$ medida a $f = \SI{1}{\kilo\hertz}$.

\begin{figure*}[t]
\centering
\begin{tikzpicture}
\begin{axis}[
    name=plot_carga,
    width=0.46\textwidth,
    height=5.9cm,
    xlabel={Corriente de Carga $I_L$ [\si{\ampere}]},
    ylabel={Tensión $V_O$ [\si{\volt}] / Potencia [\si{\watt}]},
    xmin=0, xmax=1.60,
    ymin=0, ymax=18,
    grid=both,
    grid style={line width=.1pt, draw=gray!25},
    major grid style={line width=.2pt, draw=gray!50},
    legend pos=north west,
    legend style={font=\scriptsize, fill=white, fill opacity=0.92, draw=gray!60},
    tick label style={font=\footnotesize},
    label style={font=\footnotesize\bfseries},
    title style={font=\footnotesize\bfseries, align=center},
    title={(a) Regulación de Carga, Limitación $I_{SC}$ y Balance de Potencia}
]
    \addplot[color=red!85!black, ultra thick] table[x=IL_A, y=Vo] {sim/regulacion_carga.txt};
    \addlegendentry{Tensión de salida $V_O$ [\si{\volt}]}

    \addplot[color=blue!80!black, thick, dashed] table[x=IL_A, y=PQ1_W] {sim/regulacion_carga.txt};
    \addlegendentry{Potencia en $Q_1$ $P_{C(Q_1)}$ [\si{\watt}]}

    \addplot[color=teal!85!black, thick, dashdotted] table[x=IL_A, y=P7805_W] {sim/regulacion_carga.txt};
    \addlegendentry{Potencia en \texttt{LM7805C} [\si{\watt}]}

    \addplot[color=orange!90!black, thick, dotted] table[x=IL_A, y=PLoad_W] {sim/regulacion_carga.txt};
    \addlegendentry{Potencia útil en carga $P_L$ [\si{\watt}]}

    \node[anchor=south east, font=\scriptsize\bfseries, text=red!80!black] at (axis cs:1.58, 5.2) {Corte $I_{SC} \approx \SI{1,35}{\ampere}$};
\end{axis}

\begin{axis}[
    name=plot_ripple,
    at={(plot_carga.outer north east)},
    anchor=outer north west,
    xshift=0.75cm,
    width=0.46\textwidth,
    height=5.9cm,
    xlabel={Tiempo $t$ [\si{\milli\second}] ($f_r = \SI{100}{\hertz}$, $V_{r,\text{in}} = \SI{2,0}{\volt_{\text{pp}}}$)},
    ylabel={Ondulación de Salida [\si{\micro\volt}] / Intermedia [\si{\milli\volt}]},
    xmin=0, xmax=20,
    ymin=-650, ymax=650,
    grid=both,
    grid style={line width=.1pt, draw=gray!25},
    major grid style={line width=.2pt, draw=gray!50},
    legend pos=north east,
    legend style={font=\scriptsize, fill=white, fill opacity=0.92, draw=gray!60},
    tick label style={font=\footnotesize},
    label style={font=\footnotesize\bfseries},
    title style={font=\footnotesize\bfseries, align=center},
    title={(b) Rechazo al Ripple ($\SI{2}{\volt_{\text{pp}}}$ en $V_{in}$ a \SI{100}{\hertz})}
]
    \addplot[color=purple!80!black, thick, dashed] table[x=t_ms, y=ac_VoSin_uV] {sim/ripple_tran.txt};
    \addlegendentry{$v_{o,\text{sin}}(t)$ Directo [\si{\micro\volt}] ($\SI{1116}{\micro\volt_{\text{pp}}}$)}

    \addplot[color=orange!85!black, thick, dashdotted] table[x=t_ms, y expr=\thisrow{ac_Vireg_mV}*10] {sim/ripple_tran.txt};
    \addlegendentry{$10 \times v_{i,\text{reg}}(t)$ [\si{\milli\volt}] ($\SI{48,2}{\milli\volt_{\text{pp}}}$)}

    \addplot[color=blue!85!black, ultra thick] table[x=t_ms, y expr=\thisrow{ac_VoCon_uV}*5] {sim/ripple_tran.txt};
    \addlegendentry{$5 \times v_{o,\text{con}}(t)$ Con Prereg. [\si{\micro\volt}] ($\SI{33,6}{\micro\volt_{\text{pp}}}$)}
\end{axis}
\end{tikzpicture}
\caption{Resultados de simulación SPICE: (a) característica de regulación de carga y distribución de potencia entre $Q_1$ y el \texttt{LM7805C}; (b) respuesta transitoria frente a un rizado de entrada de \SI{2,0}{\volt_{\text{pp}}} (\SI{100}{\hertz}) mostrando la reducción del ripple de salida de \SI{1116,2}{\micro\volt_{\text{pp}}} a \SI{33,6}{\micro\volt_{\text{pp}}}.}
\label{fig:simulacion_carga_ripple}
\end{figure*}

\subsubsection{Coeficiente de Temperatura ($C.T.$)}
Conforme a la Ecuación 6-44 de \cite{gonzalez_p2}, el coeficiente de temperatura expresa el cambio porcentual de la tensión de salida por grado centígrado de variación térmica (en $[\si{\percent\per\celsius}]$):
\begin{equation}
    C.T. = \pm \frac{V_{L(T_{\max})} - V_{L(T_{\min})}}{\left(T_{\max} - T_{\min}\right) V_{L(\SI{25}{\celsius})}} \times 100 = \frac{\Delta V_O / \Delta T}{V_{L(\SI{25}{\celsius})}} \times 100
    \label{eq:ct_def}
\end{equation}
Dado que el lazo de control del \texttt{LM7805C} presenta un rechazo de línea de \SI{78}{\deci\bel}, la deriva térmica del diodo Zener de la etapa prereguladora ($\alpha_{VZ} = \SI{+0,075}{\percent\per\celsius} \implies \SI{+7,5}{\milli\volt\per\celsius}$ compensada parcialmente por la deriva de $V_{BE(Q_1)} \approx \SI{-2,0}{\milli\volt\per\celsius}$, resultando $\Delta V_{i,\text{reg}}/\Delta T \approx \SI{+5,5}{\milli\volt\per\celsius}$) es atenuada en casi cuatro órdenes de magnitud por el integrado ($\approx \SI{+0,7}{\micro\volt\per\celsius}$ en la salida). En consecuencia, el $C.T.$ global queda gobernado exclusivamente por la referencia interna del \texttt{LM7805C} ($\Delta V_O / \Delta T = \SI{-1,1}{\milli\volt\per\celsius}$, Tabla~\ref{tab:datasheet_7805}):
\begin{equation}
    C.T. = \frac{\SI{-1,1}{\milli\volt\per\celsius}}{\SI{5000}{\milli\volt}} \times 100 = \mathbf{\SI{-0,022}{\percent\per\celsius}}
    \label{eq:ct_calc}
\end{equation}

\subsubsection{Tensión de Ripple ($V_{r,\text{pp}}$, $V_{\text{ef,rip}}$) y Rechazo de Ripple ($RR$)}
Si la fuente primaria de \SI{25}{\volt} proviene de un rectificador de onda completa con filtro capacitivo ($f_r = \SI{100}{\hertz}$ / \SI{120}{\hertz}) con una ondulación pico a pico de prueba $V_{r,\text{in(pp)}} = \SI{2,00}{\volt_{\text{pp}}}$, la etapa prereguladora produce una primera atenuación determinada por la Ecuación 6-76 de \cite{gonzalez_p2} (con $Z_{ZT} \parallel r_{\pi(Q_1)} = \SI{4,45}{\ohm}$):
\begin{equation}
    V_{r,\text{ireg(pp)}} \approx \SI{2,00}{\volt_{\text{pp}}} \times \frac{\SI{4,45}{\ohm}}{\SI{180}{\ohm} + \SI{4,45}{\ohm}} = \SI{48,25}{\milli\volt_{\text{pp}}}
    \label{eq:ripple_prereg}
\end{equation}
lo que equivale a un rechazo de la etapa prereguladora de $RR_{\text{prereg}} = 20\log_{10}(\SI{2000}{\milli\volt} / \SI{48,25}{\milli\volt}) = \SI{32,35}{\deci\bel}$ (y \SI{39,8}{\deci\bel} con $R_Z = \SI{430}{\ohm}$ en el \texttt{Circuito B}). A continuación, el regulador \texttt{LM7805C} operando dentro de su rango óptimo $[\SI{8}{\volt}, \SI{12}{\volt}]$ aporta su rechazo al rizado $RR_{\text{7805}} = \SI{78}{\deci\bel}$ típico (\SI{62}{\deci\bel} mínimo). En la simulación transitoria de la Figura~\ref{fig:simulacion_carga_ripple}(b) se obtiene en la salida un rizado pico a pico de tan sólo \textbf{$V_{r,\text{out(pp)}} = \SI{33,6}{\micro\volt_{\text{pp}}}$} (frente a \SI{1116,2}{\micro\volt_{\text{pp}}} sin preregulador, $RR_{\text{total}} = \SI{95,5}{\deci\bel}$). Aplicando las Ecuaciones 6-46 y 6-47 de \cite{gonzalez_p2}, el valor eficaz del ripple de salida y el factor de ripple porcentual son:
\begin{align}
    V_{\text{ef,rip}} &= \frac{V_{r,\text{out(pp)}}}{2\sqrt{3}} = \frac{\SI{33,6}{\micro\volt}}{3,464} = \mathbf{\SI{9,70}{\micro\volt_{\text{ef}}}} \label{eq:vef_rip}\\
    F.R.\ [\%] &= \frac{V_{\text{ef,rip}}}{V_L} \times 100 = \frac{\SI{9,70e-6}{\volt}}{\SI{5,00}{\volt}} \times 100 = \mathbf{\SI{0,000194}{\percent}} \label{eq:fr_pct}
\end{align}

\subsubsection{Rendimiento ($\eta$) y Distribución de Potencias}
De acuerdo con las Ecuaciones 6-1, 6-68 y 6-69 de \cite{gonzalez_p2}, el rendimiento energético de la fuente (en $[\%]$) se define como la relación porcentual entre la potencia continua entregada a la carga ($P_L = V_O \cdot I_L$) y la potencia total suministrada por la fuente primaria ($P_i = V_{in} \cdot I_{in}$):
\begin{equation}
    \eta = \frac{P_L}{P_i} \times 100 = \frac{V_O \cdot I_L}{V_{in} \left(I_L + I_d + I_{R_Z}\right)} \times 100 \approx \frac{V_O}{V_{in}} \times 100
    \label{eq:rendimiento}
\end{equation}
\begin{itemize}
    \item \textbf{En $V_{in} = \SI{25,0}{\volt}$ ($I_L = \SI{500}{\milli\ampere}$, $P_L = \SI{2,50}{\watt}$):} En el \texttt{Circuito B} ($I_{R_Z} = \SI{34,8}{\milli\ampere}$), la potencia de entrada es $P_i = \SI{25,0}{\volt} \times (\SI{0,5042}{\ampere} + \SI{0,0348}{\ampere}) = \SI{13,48}{\watt}$, resultando $\mathbf{\eta_{(\SI{25}{\volt})} = \SI{18,55}{\percent}}$ (\SI{17,10}{\percent} en el \texttt{Circuito A} y \SI{19,20}{\percent} a plena carga de \SI{1,0}{\ampere}, próximo al límite teórico ideal de un regulador lineal $\eta_{\max} = \SI{5}{\volt}/\SI{25}{\volt} = \SI{20,0}{\percent}$).
    \item \textbf{En $V_{in} = \SI{10,0}{\volt}$ ($I_L = \SI{500}{\milli\ampere}$, $P_L = \SI{2,50}{\watt}$):} La potencia de entrada se reduce a $P_i = \SI{10,0}{\volt} \times \SI{0,511}{\ampere} = \SI{5,11}{\watt}$, elevando el rendimiento a $\mathbf{\eta_{(\SI{10}{\volt})} = \SI{48,92}{\percent}}$ (límite ideal $\SI{5}{\volt}/\SI{10}{\volt} = \SI{50,0}{\percent}$).
\end{itemize}
Cabe destacar una conclusión física fundamental: en un sistema puramente lineal serie, la adición del preregulador $Q_1$ mantiene constante la caída total $(V_{in} - V_O = \SI{20}{\volt})$, pero \textbf{reparte estratégicamente la potencia disipada}: el transistor externo \texttt{TIP31C} absorbe \SI{7,79}{\watt} (\SI{77,2}{\percent}) mientras que el circuito integrado \texttt{LM7805C} sólo disipa \SI{2,31}{\watt} (\SI{22,8}{\percent}), evitando la activación de las protecciones SOA y térmica del regulador monolítico.

\subsection{Síntesis de Parámetros de Calidad y Conclusiones}
En la Tabla~\ref{tab:resumen_calidad} se compendia la totalidad de los parámetros de calidad obtenidos para la fuente diseñada, contrastándolos con la configuración sin etapa prereguladora.

\begin{table}[htbp]
\centering
\caption{ Cuadro resumen de todos los parámetros de calidad de la fuente de alimentación diseñada ($I_{L,\text{nom}} = \SI{500}{\milli\ampere}$).}
\label{tab:resumen_calidad}
\resizebox{\columnwidth}{!}{%
\begin{tabular}{llccc}
\toprule
\textbf{Parámetro de Calidad} & \textbf{Expresión / Condición} & \textbf{Sin Preregulador} & \textbf{Con Preregulador} & \textbf{Unidad} \\
\midrule
Tensión nominal de salida ($V_O$) & $V_{in} = \SI{25}{\volt}, R_L = \SI{10}{\ohm}$ & \SI{5,0058}{} & \textbf{\SI{4,9987}{}} & \si{\volt} \\
Tensión entrada al CI ($V_{i,\text{reg}}$) & Para $\SI{10}{\volt} \le V_{in} \le \SI{25}{\volt}$ & \SI{10,0}{} a \SI{25,0}{} & \textbf{\SI{8,06}{} a \SI{9,54}{}} ($\in [8, 12]$) & \si{\volt} \\
Regulación de línea ($\Delta V_O$) & $V_{in}: \SI{10}{\volt} \to \SI{25}{\volt}$ ($\SI{12}{\volt} \to \SI{25}{\volt}$) & \SI{6,79}{} (\SI{5,87}{}) & \textbf{\SI{1,70}{} (\SI{0,28}{})} & \si{\milli\volt} \\
Regulación de red ($R_{\text{red}}$) & Ec.~\eqref{eq:rred_def}, $\SI{10}{\volt}\text{--}\SI{25}{\volt}$ ($\SI{12}{\volt}\text{--}\SI{25}{\volt}$) & \SI{0,1358}{} (\SI{0,1174}{}) & \textbf{\SI{0,0340}{} (\SI{0,0055}{})} & \si{\percent} \\
Regulación de línea ($\pm \SI{10}{\percent}$ en \SI{25}{\volt}) & $V_{in}: \SI{22,5}{\volt} \to \SI{27,5}{\volt}$ & \SI{2,792}{} (\SI{0,0558}{\percent}) & \textbf{\SI{0,084}{}} (\textbf{\SI{0,0017}{\percent}}) & \si{\milli\volt} \\
Regulación de carga ($\Delta V_O$) & $I_L: \SI{5}{\milli\ampere} \to \SI{0,5}{\ampere}$ ($\to \SI{1,0}{\ampere}$) & \SI{8,51}{} (\SI{17,15}{}) & \textbf{\SI{8,51}{} (\SI{17,15}{})} & \si{\milli\volt} \\
Regulación de carga ($R_{\text{car}}$) & Ec.~\eqref{eq:rcar_def}, a \SI{0,5}{\ampere} (a \SI{1,0}{\ampere}) & \SI{0,170}{} (\SI{0,344}{}) & \textbf{\SI{0,170}{} (\SI{0,344}{})} & \si{\percent} \\
Resistencia de salida ($R_{\text{sal}}$) & Ec.~\eqref{eq:rsal_calc} / Dinámica \SI{1}{\kilo\hertz} & \SI{17,2}{} / \SI{17,0}{} & \textbf{\SI{17,2}{} / \SI{17,0}{}} & \si{\milli\ohm} \\
Coeficiente de temperatura ($C.T.$) & Ec.~\eqref{eq:ct_def}, $T_J: \SI{0}{\celsius} \to \SI{125}{\celsius}$ & \SI{-0,022}{} & \textbf{\SI{-0,022}{}} ($\Delta T_J$ \SI{77}{\percent} menor) & \si{\percent\per\celsius} \\
Rechazo de ripple ($RR_{\text{total}}$) & $f_r = \SI{100}{\hertz}\text{--}\SI{120}{\hertz}$ & \SI{65,1}{} (\SI{78}{\deci\bel} típ.) & \textbf{\SI{95,5}{}} ($\SI{32,4}{\deci\bel} + \SI{63,1}{\deci\bel}$) & \si{\deci\bel} \\
Tensión de ripple pico a pico ($V_{r,\text{pp}}$) & Para $V_{r,\text{in}} = \SI{2,0}{\volt_{\text{pp}}}$ & \SI{1116,2}{} & \textbf{\SI{33,6}{}} & \si{\micro\volt_{\text{pp}}} \\
Tensión eficaz de ripple ($V_{\text{ef,rip}}$) & Ec.~\eqref{eq:vef_rip} ($V_{r,\text{pp}} / 3,464$) & \SI{322,2}{} & \textbf{\SI{9,70}{}} & \si{\micro\volt_{\text{ef}}} \\
Potencia disipada en \texttt{LM7805C} & $V_{in} = \SI{25}{\volt}, I_L = \SI{500}{\milli\ampere}$ & \SI{10,11}{} & \textbf{\SI{2,31}{}} (\texttt{B}: \SI{2,20}{\watt}) & \si{\watt} \\
Potencia disipada en $Q_1$ (\texttt{TIP31C}) & $V_{in} = \SI{25}{\volt}, I_L = \SI{500}{\milli\ampere}$ & --- & \textbf{\SI{7,79}{}} (\texttt{B}: \SI{7,90}{\watt}) & \si{\watt} \\
Rendimiento energético ($\eta$) & Ec.~\eqref{eq:rendimiento}, en \SI{25}{\volt} / en \SI{10}{\volt} & \SI{19,67}{} / \SI{49,17}{} & \textbf{\SI{18,55}{} / \SI{48,92}{}} & \si{\percent} \\
\bottomrule
\end{tabular}%
}
\end{table}

\textbf{Conclusiones finales:}
\begin{enumerate}
    \item La incorporación de la etapa prereguladora serie discreta (\texttt{TIP31C} + referencia Zener de \SI{10}{\volt}) cumple rigurosamente con la condición de diseño de mantener la entrada del \texttt{LM7805C} dentro del intervalo óptimo $V_{i,\text{reg}} \in [\SI{8,0}{\volt}, \SI{12,0}{\volt}]$ ($\SI{8,06}{\volt}$ a \SI{9,54}{\volt}) frente a variaciones de la fuente primaria desde \SI{10}{\volt} hasta \SI{25}{\volt}.
    \item La simulación SPICE confirma que la regulación de línea mejora en un factor de $21\times$ en el rango de \SI{12}{\volt} a \SI{25}{\volt} y de $33\times$ a $60\times$ frente a fluctuaciones de red del $\pm \SI{10}{\percent}$ en torno a \SI{25}{\volt}, al tiempo que el rechazo al ripple se incrementa en más de \SI{30}{\deci\bel} (reduciendo la ondulación eficaz de salida a \SI{9,70}{\micro\volt_{\text{ef}}}) y la disipación térmica sobre el regulador integrado disminuye en un \SI{77,2}{\percent}.
    \item El transistor seleccionado (\texttt{TIP31C}) opera holgadamente dentro de su Área de Operación Segura (SOA) tanto en régimen nominal (\SI{500}{\milli\ampere}, \SI{7,90}{\watt}) como a plena capacidad (\SI{1,0}{\ampere}, \SI{15,75}{\watt}), garantizándose una temperatura de juntura segura ($T_J \le \SI{100,2}{\celsius} < T_{J,\max} = \SI{150}{\celsius}$) mediante el empleo de un disipador de aluminio extruido de $R_{th(h-a)} \le \SI{3,2}{\celsius\per\watt}$.
\end{enumerate}
